% !TEX root = ../main.tex
\chapter{Custom Function API}
\label{ch:custom_functions}

Not every operation a user wants to write fits the composite-of-
primitives pattern (\chref{ch:composite_ops}). Some operations
need a custom backward that differs from what autograd would
derive automatically; some need to interact with a non-\Lucid\
library; some need to inject a specific numerical-stability
trick. The \emph{Custom Function API} is the escape hatch.

This chapter describes the \code{lucid.autograd.Function} class:
how to define a custom op, how the autograd engine integrates it,
when to use it versus a composite, and the design considerations
that come up in practice.

The API is parity-compatible with \refframework's
\code{torch.autograd.Function}; users familiar with that interface
will see only minor naming differences.

\section{Overview}
\label{sec:custom_functions:overview}

A \code{Function} subclass defines two static methods,
\code{forward} and \code{backward}, and instructs the autograd
engine to use them when the user calls
\code{MyFunction.apply(args)}.

\begin{lstlisting}[style=lucid,language=LucidPython]
import lucid

class MyOp(lucid.autograd.Function):
    @staticmethod
    def forward(ctx, x, scale):
        ctx.save_for_backward(x)
        ctx.scale = scale
        return scale * x.exp()

    @staticmethod
    def backward(ctx, grad_out):
        (x,) = ctx.saved_tensors
        scale = ctx.scale
        grad_x = grad_out * scale * x.exp()
        return grad_x, None  # None for non-tensor `scale`

# Use it
y = MyOp.apply(x, 3.0)
y.sum().backward()
\end{lstlisting}

The \code{ctx} (context) object passes state between forward and
backward: saved tensors via \code{ctx.save\_for\_backward(...)},
non-tensor state via attribute assignment.

\section{The Function Class}
\label{sec:custom_functions:function_class}

\code{lucid.autograd.Function} is the abstract base from which
custom ops inherit. The implementation lives in
\code{lucid/autograd/function.py}; it bridges to
\code{lucid/\_C/autograd/CustomFunction.cpp} for the engine
integration.

\subsection{Required methods}
\label{ssec:custom_functions:required}

A subclass must define:

\begin{itemize}
  \item \code{forward(ctx, *inputs)}: receives the inputs (tensors
    or non-tensor arguments), returns one or more output tensors.
  \item \code{backward(ctx, *grad\_outputs)}: receives one
    cotangent per output, returns one cotangent per input
    (\code{None} for non-tensor inputs that don't need a
    gradient).
\end{itemize}

Both are \code{@staticmethod} decorated. The class itself is never
instantiated; \code{MyOp.apply(args)} is the call.

\subsection{The Context object}
\label{ssec:custom_functions:context}

\code{ctx} is a per-call object that lives from the forward
through the backward. The user uses it to:

\begin{itemize}
  \item Save tensors needed by backward:
    \code{ctx.save\_for\_backward(t1, t2, ...)}. Later retrieved
    in backward as \code{ctx.saved\_tensors}.
  \item Stash non-tensor state: \code{ctx.my\_attr = value} in
    forward, \code{ctx.my\_attr} in backward.
  \item Inspect autograd state: \code{ctx.needs\_input\_grad} is
    a tuple of booleans indicating which inputs require grad.
  \item Mark inputs as not requiring grad even though they
    nominally do: \code{ctx.mark\_non\_differentiable(t)}.
\end{itemize}

\subsection{The apply method}
\label{ssec:custom_functions:apply}

\code{MyOp.apply(*args)} is the user-facing entry point. Under
the hood:

\begin{enumerate}
  \item Construct a fresh \code{ctx} object.
  \item Call \code{forward(ctx, *args)} to get the outputs.
  \item If any input requires grad, construct a
    \code{CustomFunctionNode} carrying \code{ctx} and the user's
    \code{backward} method. Wire it as the outputs'
    \code{grad\_fn}.
  \item Return the outputs.
\end{enumerate}

The \code{CustomFunctionNode} is a regular autograd \code{Node}
(\chref{ch:autograd_engine}); the engine treats it like any
other node during backward traversal.

\section{Worked Example: A Custom Activation}
\label{sec:custom_functions:custom_activation}

Suppose we want to implement a non-standard activation,
\code{swish\_beta(x) = x * sigmoid(beta * x)} with a learnable
$\beta$. As a composite, the gradient through $\beta$ would be
constructed by autograd automatically, but for didactic purposes
we write it as a Function with an explicit backward.

\begin{lstlisting}[style=lucid,language=LucidPython]
import lucid
from lucid.autograd import Function

class SwishBeta(Function):
    @staticmethod
    def forward(ctx, x, beta):
        sig = (beta * x).sigmoid()
        out = x * sig
        ctx.save_for_backward(x, beta, sig, out)
        return out

    @staticmethod
    def backward(ctx, grad_out):
        x, beta, sig, out = ctx.saved_tensors

        # d/dx [x * sigmoid(beta*x)]
        #   = sigmoid(beta*x) + x * beta * sigmoid(beta*x) * (1 - sigmoid(beta*x))
        d_sig = sig * (1 - sig)
        grad_x = grad_out * (sig + x * beta * d_sig)

        # d/dbeta [x * sigmoid(beta*x)]
        #   = x * x * sigmoid(beta*x) * (1 - sigmoid(beta*x))
        grad_beta_per = grad_out * x * x * d_sig
        grad_beta = grad_beta_per.sum()  # reduce to scalar

        return grad_x, grad_beta
\end{lstlisting}

The forward saves four tensors (input, beta, the sigmoid
intermediate, and the output). The backward derives the
gradients in closed form using the saved quantities.

\subsection{Why save the sigmoid intermediate}
\label{ssec:custom_functions:save_intermediate}

The sigmoid intermediate could be recomputed in backward from
\code{x} and \code{beta}, but saving it costs $O(N)$ memory and
saves one transcendental kernel evaluation. The trade-off favours
saving for medium-to-large $x$. For very large $x$ where memory
matters more than compute, the user can choose not to save and
recompute instead --- \code{gradient\_checkpointing}
(\chref{ch:gradient_checkpointing}) generalises this trade-off.

\subsection{When this would be a primitive}
\label{ssec:custom_functions:could_be_prim}

If \code{SwishBeta} were used in every layer of a deep network,
the per-call framework overhead and the saved-tensor cost would
add up. At that point we would promote it to a primitive
(\chref{ch:op_registry}), with a C++ kernel for the forward and
a registered \code{BackwardSpec}. The Custom Function path is for
operations not common enough to justify the primitive build cost.

\section{Worked Example: Wrapping an External Library}
\label{sec:custom_functions:external}

Custom Function is the right tool when the operation is provided
by an external library (a third-party C kernel, a MPSGraph
operation \Lucid\ does not yet wrap, etc.).

\begin{lstlisting}[style=lucid,language=LucidPython]
class ExternalGEMM(Function):
    @staticmethod
    def forward(ctx, A, B):
        # Call an external C library
        out = _external_gemm(A, B)
        # Wrap result in a Lucid tensor sharing the bytes
        out_tensor = lucid.from_dlpack(out)
        ctx.save_for_backward(A, B)
        return out_tensor

    @staticmethod
    def backward(ctx, grad_out):
        A, B = ctx.saved_tensors
        # We know the gradient even though the forward was external
        grad_A = grad_out @ B.transpose(-1, -2)
        grad_B = A.transpose(-1, -2) @ grad_out
        return grad_A, grad_B
\end{lstlisting}

The pattern: the forward routes through the external call, but
the backward stays in \Lucid\ primitives. The framework only sees
a single \code{ExternalGEMM} node in the backward graph; the
external call's internals are opaque.

\section{The needs\_input\_grad Optimisation}
\label{sec:custom_functions:needs_input_grad}

\code{ctx.needs\_input\_grad} is a tuple of booleans matching the
input order. The user can inspect it in \code{backward} to skip
computing gradients that the autograd engine will discard anyway.

\begin{lstlisting}[style=lucid,language=LucidPython]
class MyOp(Function):
    @staticmethod
    def forward(ctx, x, y, z):
        ctx.save_for_backward(x, y, z)
        return some_op(x, y, z)

    @staticmethod
    def backward(ctx, grad_out):
        x, y, z = ctx.saved_tensors
        grad_x = grad_y = grad_z = None
        if ctx.needs_input_grad[0]:
            grad_x = ...  # expensive computation
        if ctx.needs_input_grad[1]:
            grad_y = ...
        if ctx.needs_input_grad[2]:
            grad_z = ...
        return grad_x, grad_y, grad_z
\end{lstlisting}

The optimisation matters when some gradients are expensive and
not always needed. For typical training where every parameter
requires grad, every \code{needs\_input\_grad[i]} is True and the
optimisation is moot.

\section{Validating Custom Functions with gradcheck}
\label{sec:custom_functions:gradcheck}

A custom backward is a place where bugs hide. \Lucid's
\code{gradcheck} (\chref{ch:gradcheck}) compares the analytical
backward against a finite-difference approximation and raises if
they disagree.

\begin{lstlisting}[style=lucid,language=LucidPython]
from lucid.autograd import gradcheck

x = lucid.randn(5, requires_grad=True, dtype=lucid.float64)
beta = lucid.tensor(0.7, requires_grad=True, dtype=lucid.float64)

assert gradcheck(SwishBeta.apply, (x, beta), eps=1e-6, atol=1e-4)
\end{lstlisting}

\subsection{Why FP64 for gradcheck}
\label{ssec:custom_functions:gradcheck_fp64}

The numerical-derivative error is $O(\sqrt{\epsilon})$ where
$\epsilon$ is machine precision. In FP32 that's $\sim 10^{-4}$,
which is the threshold at which most gradchecks barely pass; in
FP64 it's $\sim 10^{-8}$, leaving room for tight tolerances.

\Lucid\ does not require FP64 for runtime use, but gradcheck
should be in FP64 to give a meaningful signal. The tolerance for
FP32 gradcheck has to be loose (\code{atol=1e-3}) and is less
useful as a debugging tool.

\section{Composing Functions}
\label{sec:custom_functions:composing}

Custom Functions compose: one Function's forward can call
\code{apply} on another. The autograd engine builds the chain
naturally.

The chain has implications for higher-order differentiation
(\chref{ch:higher_order}): a Function's backward that itself
includes a Function call participates in the second-order
gradient.

\section{When Not to Use Custom Function}
\label{sec:custom_functions:when_not}

The decision tree:

\begin{enumerate}
  \item Can the op be expressed as a composition of existing
    primitives? Use a composite
    (\chref{ch:composite_ops}). The autograd-derived backward is
    correct by construction.
  \item Is the op called millions of times per training step on
    a hot path? Build a primitive
    (\chref{ch:op_registry}).
  \item Otherwise: Custom Function.
\end{enumerate}

The middle ground (custom op needed, not on hot path, not
expressible as composite) is the Custom Function use case. In
practice it covers wrappers around external libraries, research
prototype ops, and any case where the auto-derived backward is
numerically problematic and needs an analytic replacement.

\section{Implementation Notes}
\label{sec:custom_functions:impl}

The Python \code{Function} class is in
\code{lucid/autograd/function.py}. It wraps the C++
\code{CustomFunctionNode} in \code{lucid/\_C/autograd/CustomFunction.cpp}.

\subsection{The Python-to-C++ boundary}
\label{ssec:custom_functions:py_cpp_boundary}

The custom \code{backward} is a Python function. The C++ engine,
when it walks the backward graph and reaches a
\code{CustomFunctionNode}, calls back into Python to invoke the
function. This requires holding the GIL during the call, which is
the one place in the engine that the GIL is acquired during
backward.

For pure-\Lucid-primitive backwards inside the custom function,
the GIL is then released for the duration of each primitive call.
The net overhead is a few microseconds per Custom Function call
beyond a plain primitive node.

\subsection{Memory model}
\label{ssec:custom_functions:memory}

The \code{ctx} object is held by the \code{CustomFunctionNode}
for the lifetime of the backward graph. Anything assigned to
\code{ctx} (saved tensors, Python attributes) lives until backward
completes. The user should avoid storing large Python objects
on \code{ctx} if memory is tight.

\subsection{Thread safety}
\label{ssec:custom_functions:thread_safety}

Custom Function's \code{forward} and \code{backward} are called
on the calling thread. Multiple concurrent backward passes (from
different model invocations) instantiate different \code{ctx}
objects and operate independently. The user's custom code is
responsible for any internal thread safety (which is rare in
practice).

\section{Worked Example: A Straight-Through Estimator}
\label{sec:custom_functions:ste_example}

A common pattern in models that need a discrete decision (binary
gating, quantisation, hard attention) is the \emph{straight-through
estimator} (STE): the forward applies a non-differentiable rounding
or thresholding, but the backward passes the gradient through
unchanged, as if the rounding were the identity. This is a
mathematically illegitimate but empirically effective approximation
that lets gradient-based training work with discrete decisions.

The STE is naturally a Custom Function:

\begin{lstlisting}[style=lucid,language=LucidPython]
class BinaryQuantize(Function):
    @staticmethod
    def forward(ctx, x):
        # Forward: round to 0 or 1 based on sign
        return (x > 0).astype(x.dtype)

    @staticmethod
    def backward(ctx, grad_out):
        # Straight-through: gradient passes unchanged
        return grad_out

binarize = BinaryQuantize.apply
\end{lstlisting}

This is two-line backward, but it could not be expressed as a
composite because no composition of differentiable primitives
produces the identity gradient through a step function. The Custom
Function escape hatch is the only clean way.

\subsection{Variants}
\label{ssec:custom_functions:ste_variants}

The STE family has several refinements:

\begin{itemize}
  \item \textbf{Clipped STE}: pass gradient through only when the
    input is in a safe range, zero otherwise.
  \item \textbf{Soft STE}: blend the analytical sub-gradient with
    the straight-through gradient.
  \item \textbf{Annealed STE}: start with a smooth surrogate
    (sigmoid) and anneal toward hard binarisation.
\end{itemize}

Each is a 5--10 line Custom Function in \Lucid.

\section{Worked Example: A Reverse-Mode Differentiable Solver}
\label{sec:custom_functions:solver_example}

Another use case: wrapping a numerical solver that does not admit
straightforward AD. Consider a function that returns the root of a
nonlinear equation via Newton's method:

\begin{lstlisting}[style=lucid,language=LucidPython]
class ImplicitRoot(Function):
    @staticmethod
    def forward(ctx, params, x0):
        # Solve g(x; params) = 0 via Newton iteration
        x = x0
        for _ in range(50):
            with lucid.no_grad():
                gx = _g(x, params)
                gx_x = _g_jac_x(x, params)
                x = x - gx / gx_x
        ctx.save_for_backward(x, params)
        return x

    @staticmethod
    def backward(ctx, grad_out):
        # Implicit function theorem:
        # dx/dparams = -(d g/d x)^{-1} (d g/d params)
        x, params = ctx.saved_tensors
        gx_x = _g_jac_x(x, params)
        gx_p = _g_jac_p(x, params)
        grad_params = -(grad_out / gx_x) * gx_p
        return grad_params, None  # x0 has no gradient
\end{lstlisting}

The forward runs an iterative solver that is not amenable to
unrolling through autograd (the iteration count is dynamic and the
intermediate states are not useful for gradients). The backward
applies the implicit function theorem: the gradient with respect
to the parameters is computed in closed form from the Jacobians of
the residual function $g$ at the solution.

This pattern --- iterative forward, analytical backward via the
implicit function theorem --- is the standard mechanism for
differentiable optimisation layers (DEQ, OptNet, and similar
modern architectures).

\section{Patterns and Anti-patterns}
\label{sec:custom_functions:patterns}

A catalogue of common mistakes when writing a Custom Function.

\subsection{Pattern: clone before in-place}
\label{ssec:custom_functions:clone_before}

If the forward modifies an input in place, the saved-for-backward
version is the modified one, and the backward gets wrong values.
The fix is to clone first:

\begin{lstlisting}[style=lucid,language=LucidPython]
@staticmethod
def forward(ctx, x):
    x_clone = x.clone()
    x_clone.add_(1.0)        # safe: modifies clone
    ctx.save_for_backward(x_clone)  # backward sees the modified value
    return x_clone
\end{lstlisting}

\subsection{Anti-pattern: saving Python objects on ctx}
\label{ssec:custom_functions:anti_pyobj}

Avoid storing large Python objects on \code{ctx}: they live for
the entire backward graph's lifetime and contribute to memory
pressure. Save only what backward needs.

\subsection{Anti-pattern: forgetting to handle non-tensor inputs}
\label{ssec:custom_functions:anti_nontensor}

If a Function takes non-tensor arguments (a scalar, a string, a
configuration dict), the backward must return \code{None} for
each. Forgetting produces a shape-mismatch error in the engine:

\begin{lstlisting}[style=lucid,language=LucidPython]
class BadFunc(Function):
    @staticmethod
    def forward(ctx, x, scale):  # scale is a Python float
        ctx.scale = scale
        return scale * x

    @staticmethod
    def backward(ctx, grad_out):
        return ctx.scale * grad_out   # WRONG: returns only one value
        # Should be:
        # return ctx.scale * grad_out, None
\end{lstlisting}

The framework expects one return value per input; \code{None} for
non-tensor inputs.

\subsection{Anti-pattern: writing to non-saved state}
\label{ssec:custom_functions:anti_state}

A Function whose forward modifies global state (a counter, a
cache, a log) is not reproducible under autograd: a re-run during
checkpointing (\chref{ch:gradient_checkpointing}) or
gradgradcheck (\chref{ch:gradcheck}) will fire the side effect
twice. Functions should be pure.

\section{Summary and Forward References}
\label{sec:custom_functions:summary}

The Custom Function API is \Lucid's extension mechanism for
operations that cannot be expressed as compositions of existing
primitives and are not common enough to justify a C++ primitive.
The user defines a \code{Function} subclass with
\code{forward(ctx, *inputs)} and \code{backward(ctx, *grad\_outputs)}
static methods; the framework integrates the custom op into the
autograd graph through a \code{CustomFunctionNode}. The
\code{ctx} object carries state between forward and backward.

\code{gradcheck} (\chref{ch:gradcheck}) is the recommended
validation step for any custom backward.

The next chapter, \chref{ch:gradient_checkpointing}, treats the
memory-saving recomputation pattern that interacts with the
saved-tensor mechanism. \chref{ch:higher_order} treats
higher-order differentiation, in which custom backwards
participate naturally.

\begin{lucidnote}
For the user: \code{Function} is the right tool for one-off custom
ops. If you are writing one and the backward is more than a few
lines, consider whether the composite path
(\chref{ch:composite_ops}) would suffice; if the backward is
extremely performance-critical, consider whether a primitive
(\chref{ch:op_registry}) is warranted.
\end{lucidnote}
